% 2-2-statistical.tex
%  Budget: 2pp.
%  Rules: every numeral in \lr{}; cite only from references.bib;
%  em-dash must be \lr{---} (bare --- is a tofu box in B Nazanin).

\section{روش‌های آماری کلاسیک}

نخستین نسل مدل‌های پیش‌بینی قیمت برق بر پایه سری‌های زمانی آماری بنا شده
بودند. اهمیت این خانواده در پژوهش حاضر صرفاً تاریخی نیست: یکی از پنج مدل
این پایان‌نامه، یعنی \lr{SARIMAX}، مستقیماً از همین سنت می‌آید، و مدل
\lr{LEAR-LASSO} نیز در ساختار خطی خود به آن نزدیک است.

نوگالس و همکاران \cite{nogales_forecasting_2002} دو ابزار پیش‌بینی مبتنی بر
مدل‌های سری زمانی برای قیمت روز-بعد ارائه می‌کنند و انگیزه کار خود را
نیاز تولیدکنندگان و مصرف‌کنندگان به پیش‌بینی دقیق برای طراحی راهبرد پیشنهاد
قیمت می‌دانند. این چارچوب انگیزشی \lr{---} پیش‌بینی به‌مثابه ورودی تصمیم
اقتصادی و نه هدفی فی‌نفسه \lr{---} در سراسر ادبیات بعدی تکرار می‌شود.

کونجو و همکاران \cite{conejo_day-ahead_2005} گامی روش‌شناختی مهم برمی‌دارند:
به‌جای برازش مستقیم یک مدل \lr{ARIMA} بر سری خام قیمت، ابتدا سری را با
تبدیل موجک به مجموعه‌ای از سری‌های خوش‌رفتارتر تجزیه می‌کنند، سپس بر هر
مؤلفه مدل \lr{ARIMA} برازش می‌دهند و در پایان نتایج را بازترکیب می‌کنند.
منطق این کار آن است که سری قیمت برق به‌طور مستقیم بدرفتار است و پیش‌پردازش
می‌تواند بخشی از دشواری را پیش از مدل‌سازی حل کند. همین ایده، بعدها در
مدل‌های ترکیبی مبتنی بر یادگیری عمیق نیز بازتولید شد (بخش \lr{2-5}).

\subsection*{چرا مدل‌های خطی هنوز کنار گذاشته نشده‌اند}

با وجود گسترش روش‌های غیرخطی، مدل‌های آماری کلاسیک همچنان در این حوزه حضور
دارند و این حضور دلایل روشنی دارد. پوجی و همکاران
\cite{poggi_electricity_2023} در مطالعه‌ای بر بازار آلمان، هم روش‌های
استنباطی آماری سنتی و هم روش‌های یادگیری عمیق جدیدتر را تحلیل می‌کنند و
راهکاری مبتنی بر تجزیه سری و استفاده از متغیرهای مجازی روزهای هفته پیشنهاد
می‌دهند. حضور همزمان هر دو خانواده در یک مطالعه واحد نشان می‌دهد که برتری
یکی بر دیگری بدیهی تلقی نمی‌شود.

گلاسی و همکاران \cite{ghelasi_day-ahead_2025} نیز بر مدل‌های اقتصادسنجی
تمرکز می‌کنند و اهمیت آن‌ها را در افق‌های میان‌مدت و بلندمدت برجسته
می‌سازند. آنان به‌صراحت به بحران انرژی اروپا در سال \lr{2022} و نوسانات
بی‌سابقه ناشی از آن اشاره می‌کنند و این نوسانات را چالشی برای مدیریت
عملیاتی و ریسک می‌دانند.

این نکته اخیر، به یکی از استدلال‌های محوری پایان‌نامه حاضر متصل می‌شود.
مدل‌های خطی، برخلاف مدل‌های مبتنی بر درخت، توانایی برون‌یابی\LTRfootnote{Extrapolation}
دارند: وقتی مقدار یک متغیر ورودی خارج از دامنه‌ای قرار گیرد که مدل بر آن
آموزش دیده است، مدل خطی همچنان مقداری متناسب تولید می‌کند، در حالی که مدل
درختی به آخرین برگ خود محدود می‌ماند. این تفاوت ساختاری، در شرایط عادی
اهمیت چندانی ندارد، اما چنان‌که در فصل پنجم نشان داده می‌شود، در آزمون
برون‌توزیع این پژوهش به عاملی تعیین‌کننده بدل شد.

بنابراین جایگاه این خانواده در کار حاضر دوگانه است: هم به‌عنوان رقیب جدی در
مقایسه اصلی، و هم به‌عنوان مرجعی برای تفسیر رفتار مدل‌ها در شرایطی که
مفروضات آموزش نقض می‌شود.
